\documentclass[UTF8,fontset=none,11pt,a4paper]{ctexart}
\usepackage{ipara-handbook}
\handbooksetup{DS-06 Union-Find 与集合划分}{408 · 数据结构 DS-06}{Partition · Representative · Find · Union}
\begin{document}
\handbooktitle{Union-Find 与集合划分}{动态合并后，怎样快速回答“是否同一集合”}{408 · 数据结构 Topic DS-06}

\begin{corebox}{母问题：集合不断合并时，怎样快速判断两个元素现在是不是同一组？}
并查集的核心过程可以写成
\[
\text{若干互不相交的集合}\to\text{每组选一个根作代表}\to\text{Find 找根 / Union 合并根}\to\text{压短查找路径}\to\text{快速判断连通}.
\]
每个集合用一棵父指针树保存，但我们真正关心的是“哪些元素属于同一集合”，而不是树画成什么形状。\code{find(x)} 沿父指针找到该集合的根；\code{union(a,b)} 只在两端根不同的时候合并这两个集合。任何操作后，每个节点沿父指针都必须最终到达唯一根，根的数量也必须等于当前集合数。
\end{corebox}
\begin{methodbox}{本册定位与 Owner 边界}
\begin{itemize}
\item \kw{Owns：}MakeSet、Find、Union、代表元、按秩/大小合并、路径压缩与动态连通语义。
\item \kw{Uses：}数组下标和成本向量；向 DS-B03/DS08 输出 `connected/unite` 接口。
\item \kw{Stops：}一般树遍历、Kruskal 的边权排序与生成树正确性由其他 Owner 负责。
\end{itemize}
\end{methodbox}
\tableofcontents\newpage

\section{每次操作后都要保证：每个元素最终都能找到唯一根}
初始时每个元素自成集合，$parent[i]=i$。父指针反复追踪必到根；根是自己的父节点。若两个元素根相同，则属于同一集合；若根不同，合并时把一个根挂到另一个根下，并将集合数减一。重复合并不改变分区。

并查集的“集合”是一个随操作序列变化的等价关系：自反、对称、传递性始终成立。\code{find} 只返回代表元，不改变等价类；路径压缩虽会改写多个 parent 字段，但改写前后的根相同，因此不改变答案。\code{union} 必须先 Find 两端再合并根，不能把任意节点直接挂到任意节点，否则可能破坏“只有根才代表集合”的不变量。

若采用一个数组同时保存父指针和大小，可令根节点存负的集合大小、非根节点存父下标：$parent[root]=-size$，合并时把较小集合的根挂到较大集合，并更新负值。这省去单独的 rank/size 数组，但比较时必须只对根读取负值，非根位置的数值没有大小意义。初始化、越界和 $n=0$ 都应作为独立边界测试。

\section{为什么还要按秩/按大小合并和路径压缩}
若只把一个根挂到另一个根下，连续合并可能生成长链，使 Find 退化到 $O(n)$。按秩合并总把较矮树挂到较高树，限制树高；路径压缩在 Find 返回时让沿途节点直接指向根，摊还复杂度达到近似常数 $O(\alpha(n))$。

\subsection{朴素合并的最坏形状与两种优化的分工}
如果每次都令 $root_a.parent=root_b$，合并序列可以把元素串成一条链：
\[
0\to1\to2\to\cdots\to(n-1).
\]
此时一次 Find 可能追踪 $n$ 条父边。按秩合并保存的是“树高的上界”或等价的规模信息：只有当两棵树秩相同时，合并后新根的秩才增加一；把低秩根挂到高秩根不会破坏上界。按大小合并使用集合元素数代替秩，证明思路相同。

路径压缩利用的是“这次 Find 已经知道根是谁”：递归回溯时，让查找路径上的节点直接指向根；迭代版也可以先记录路径，再统一改父指针。它不改变集合成员，只改变以后查找时要经过多少层。两种优化分工不同：按秩/按大小合并尽量不让新树长高；路径压缩则把已经走过的长路径压短。只使用其中一种也不会影响正确性，但效率上界不同。

按秩合并的高度证明是对数级：只有两棵相同秩的树合并时秩才加一，而秩为 $r$ 的根至少包含 $2^r$ 个节点，因此 $r\le\lfloor\log_2 n\rfloor$。仅有按秩（不压缩）时单次 Find 为 $O(\log n)$；仅有路径压缩（不按秩）时可在特定序列下变差；两者同时使用时，任意 $m$ 次操作总时间为 $O(m\,\alpha(n))$，其中反阿克曼函数在实际输入范围内极小，但它不是严格的常数承诺。若题目改问“随机合并后的平均/最大树高”，这已经不是并查集 ADT 自带的结论，而是\kw{挂接规则 + 是否压缩 + 随机模型 + 操作次数}共同决定的概率问题；这些条件没有给全时，不能把某个 $O(\log n)$ 或 $O(n)$ 结论当成无条件定理。

\subsection{代表元只是这组元素当前选出的根，不是固定身份}
代表元用于回答“两个元素是否同集合”，合并之后完全可能改变。题目若额外要求某个特定元素始终作为代表元，就必须另外规定合并策略；普通按秩合并并不保证“编号小的元素一定当根”。并查集也不会自动保存集合内部顺序、两点路径长度或完整成员列表；如果题目需要这些信息，就要额外维护，不能只从 parent 数组中硬推。

\subsection{用在无向图里：先判断两端是否同根，再决定是否合并}
把无向边按某种顺序处理时，若边两端 Find 结果不同，加入该边不会在当前森林中形成环，并执行 Union；若代表元相同，则该边会连接同一连通分量内部，应拒绝。Kruskal 还需要按边权排序、停止于 $n-1$ 条边并证明得到最小生成树，这些由 DS08 Own；DS06 只保证每次“同组判定—合并”状态转移正确。

若题目只问连通分量数量，初始为 $n$ 个集合，每次成功合并减一，重复边不减。若题目给自环或重复边，Find 会返回同一代表元，正好把这些边识别为不改变分区的事件。不要把“边数达到 $n-1$”当成无条件连通；只有成功合并次数达到 $n-1$ 且不存在被拒绝的分支时，才能推出一棵生成树。

并查集适合“只增不删”的无向连通性：边一旦合并就不能直接撤销，不能回答带时间顺序的删除边问题，也不能返回两点之间的具体路径、距离或最短路。若问题要求这些信息，应改用 BFS/DFS、动态树、离线回滚或其他结构，并明确它们的 Owner。并查集也不区分有向边方向；把有向图的可达性直接交给 Union-Find 会把方向信息丢失。

\section{代码：代表元、路径压缩与按秩合并}
完整实现位于 \codepath{code/ds06_union_find.hpp}，测试位于 \codepath{code/ds06_union_find_test.cpp}；正文直接引用真实源码。
\lstinputlisting[language=C++,firstline=11,lastline=58,firstnumber=1,numbers=left,caption={Find 的路径压缩与 Union 的按秩合并}]{30_408/10_数据结构/06_UnionFind与集合划分/code/ds06_union_find.hpp}
`find` 的递归返回值就是代表元；回溯赋值是路径压缩。`unite` 先分别 Find，再比较秩，最后只减少一次 `components`。根相同的分支必须提前返回，否则重复合并会错误减少集合数。
\lstinputlisting[language=C++,firstline=60,lastline=82,firstnumber=1,numbers=left,caption={分区不变量和越界边界}]{30_408/10_数据结构/06_UnionFind与集合划分/code/ds06_union_find.hpp}

\subsection{测试：重复合并、压缩与空全集}
\lstinputlisting[language=C++,firstline=7,lastline=52,firstnumber=1,numbers=left,caption={代表元一致性和集合数变化测试}]{30_408/10_数据结构/06_UnionFind与集合划分/code/ds06_union_find_test.cpp}
\begin{lstlisting}[language=bash]
clang++ -std=c++17 -Wall -Wextra -Werror -pedantic \\
  code/ds06_union_find_test.cpp -o /tmp/ds06_test
/tmp/ds06_test
\end{lstlisting}

\section{复杂度与边界}
\begin{center}\small
\begin{tabularx}{\linewidth}{L{2.4cm}C{2cm}YY}\toprule
操作 & 摊还复杂度 & 保证 & 边界\\\midrule
Find/Union & $O(\alpha(n))$ & 路径压缩 + 按秩/大小 & 空集合、越界\\
MakeSet & $O(n)$ & 每项自为根 & $n=0$ 合法\\
空间 & $O(n)$ & parent/rank 两数组 & 不保存集合内部顺序\\
\bottomrule\end{tabularx}\end{center}
\begin{warnbox}{不要把父指针树当成题目真正研究的树}
并查集的父指针树只是为了找代表元。它不负责按层遍历、输出两点路径，也不保持关键字有序。若把这棵实现树当作普通树去研究层次、路径或节点顺序，就会把“怎么存集合”误当成“题目真正要问的结构”。
\end{warnbox}

\section{遇到并查集题时按什么顺序做}
先找当前每个集合的根；合并前只比较两个根是否相同；同根则这次 Union 不改变集合数，不同根才真正合并。若题目问复杂度，再检查是否同时使用了路径压缩和按秩/按大小合并。忘掉代码时只需重新问：\kw{根怎么表示？Find 沿什么指针找到根？Union 把哪个根挂到哪个根？哪些节点会被路径压缩？集合数什么时候减一？}
\begin{corebox}{跨册交接}
DS06 向 DS-B03 和 DS08 只提供 `connected/unite`。Kruskal 的排序、选边和无环性证明由 DS08 Own；本册只保证每次加入边前后动态分区不变量。
\end{corebox}
旧总册关于代表元、路径压缩、按秩合并和 Kruskal 接口的内容已吸收；带权并查集、可回滚并查集和持久化版本作为 Extension。
\end{document}
